\section{A sound use of an uncleared counter: transfer the residual into the parameter}
\label{sec:dirty-counter-salt}

This section distinguishes three different proposals.  Replacing the four
equality tests by a comparator which leaves either scratch register dirty
is unsafe.  Leaving only register~6 dirty avoids a forged cache, but it
still corrupts the first freshness test and the proof enumeration.  A
repaired variant has an ordinary, unconditional correctness argument:
clear register~6 between the two freshness tests, and transfer its value
into the additive parameter register~4 at the start of every command.
The latter transfer deliberately changes the enumerated commands' effective
parameters.  The resulting proof search is nevertheless sound and complete.
The measured repaired source has 208 DAG nodes; this is a semantic result,
not a claim of a machine with at most 218 states.

\subsection{The asymmetric comparator and its exact residual}

Write $D_i$ for decrement, $I_i$ for increment, $C_j$ for continue through
$j$ right ancestors, and $B_j$ for break through $j$ left ancestors, with
the same source-tree conventions as the K219 report.  The asymmetric
comparator found in the independent encoding search is
\[
 E_a=\bigl((D_5,((D_6,C_2),(I_6,(D_5,C_1)))),(D_6,B_1)\bigr).
\]
For inputs $(r_5,r_6)=(u,v)$ it terminates, preserves every other register,
and has the following contract:
\[
\begin{array}{c|c|c}
\text{condition}&\text{outcome}&(r_5,r_6)\\\hline
u=v&\mathrm{Normal}&(0,0)\\
u>v&\mathrm{Break}_1&(0,0)\\
u<v&\mathrm{Break}_1&(0,v-u-1).
\end{array}
\]
To prove the contract, cancel one unit from each register while both are
positive.  If register~6 becomes zero first, the $I_6$ branch installs one
unit in register~6 and clears register~5; the final $D_6$ consumes that
unit and $B_1$ reports inequality.  If register~5 becomes zero first,
the final $D_6$ either reports equality at zero or leaves the stated
positive-side residual.  Induction on $u+v$ proves termination.

At the NIP test the second input is exactly~1.  Hence
$\max(1-u-1,0)=0$, and this comparator is an exact clean replacement at
that call site for every cache value.

The two AX tests compare $(b,c)$ and $(a,c)$ in that order.  Without a
repair, the second pair of inputs becomes
\[
 (a,\ c+\max(c-b-1,0)).
\]
For example, $(a,b,c)=(2,0,2)$ first leaves a residual~1, then compares
$(2,3)$.  Both comparisons report inequality and the final scratch is
clean, so the unmodified caller admits the prohibited bundle
$\operatorname{bundle}(2,0,2)$.  This is a concrete incorrect gate,
not merely a missing scratch invariant.

For comparison, the more aggressive comparator that may leave register~5
dirty has a second, direct defect.  Start the AX gate with
$(a,b,c)=(3,0,1)$ and clean scratch registers.  Its first comparison
$(0,1)$ leaves $(r_5,r_6)=(0,0)$, and its second comparison $(3,1)$ leaves
$(r_5,r_6)=(1,0)$.  Register~5 is the cache tested by NIP.  An immediately
following NIP therefore halts without a cached proof of code~1.  This is
an exact source-command counterexample to the inherited soundness
invariant.  It does not by itself prove that the particular command
sequence is reached by that variant from its blank-tape initialization.

\subsection{Why leaving the residual in the input reader is insufficient}

Let $T(d)=d(d+1)/2$ and use the existing Cantor pairing
$P(a,b)=T(a+b)+b$.  The source unpairing loop, entered on a stream $n$
with register~6 equal to $d$ and register~7 zero, returns
\[
 \operatorname{unP}(n+T(d)).
\]
Indeed its internal pair starts at $(d,0)$, whose Cantor code is $T(d)$,
and each consumed stream unit advances that code by one.  Its returned
tail is still at most $n$: that tail starts at zero, increases by at most
one per unit, and is occasionally reset to zero.

Thus, while the proof stream is nonempty, the residual merely changes
which natural command fields are read.  At an empty stream, however, the
counter-reloading loop copies the next-proof counter $N$ into register~6
additively.  It reads candidate $N$ but increments the next-proof counter
to $N+d+1$.  Candidate numbers can therefore be skipped.  The original
fairness theorem does not apply to this proposal.

\subsection{The repaired enumerator}

Make exactly these source changes:
\begin{enumerate}
\item Use $E_a$ at all four equality call sites.
\item Clear register~6 between the two AX comparisons, before copying the
saved $c$ into register~6 for the second comparison.
\item At the root of every command iteration, execute the existing
transfer $6\mathbin{\to}4$, namely $(D_6,(I_4,C_2))$.
\end{enumerate}
The third operation preserves the proof stream, next-proof counter,
cache, current formula, and formula stack.  It only adds the residual
to the free, additive parameter.  Consequently every command is decoded
by the original clean reader and the next-proof counter increments by
exactly one.

At boundaries immediately after this transfer, the altered abstract
machine has the original tuple $(t,w,d,h)$.  Its only semantic changes are
the following updates to $d$:
\[
\begin{array}{c|c}
\text{command case}&\text{new parameter value}\\\hline
\mathrm{AX}(c,x,y,u),\ a=d+x
 &\max(c-a-1,0)\\
\text{failed }\mathrm{MP}(x),\ a=d+x
 &a+\max(I(t,a)-h-1,0).
\end{array}
\]
Successful MP, NIP, and DROP have their original behavior.  AX still
emits exactly the original logical rows and the exactly gated original
bundle, and clears its cache.  Failed MP still emits truth code~0 and
clears its cache.  Since $d$ is an arbitrary natural parameter, induction
on commands proves the original theorem-valued stack/cache invariant.
NIP therefore halts only on a derivation of code~1.

\subsection{An explicit completeness repair}

The following proof-compiler changes establish completeness; they do not
need any new source instructions.  Write $\mathrm{AX}_0$ for the command
with $c=x=y=u=0$.  When entered with cache zero, it has $b=c=0$, so the
bundle is suppressed, its final register~6 residual is zero, and its
parameter and cache are both zero.  As in K219, it pushes the old current
formula, one generalization, and eight logical rows.

First replace the compiler's DISCARD block by
\[
 \mathrm{DISCARD}_a
   =[\mathrm{MP}(1),\mathrm{AX}_0]\;\mathbin{+\!+}\;
     \mathrm{NIP}^{10}.
\]
It maps $(t,w,d,0)$ to $(0,w,0,0)$ for every natural $d$.
The first MP fails because $I(t,d+1)>0$, whereas the cache is zero.
Its arbitrary residual is absorbed into $d$ before $\mathrm{AX}_0$;
that command then resets the parameter.  The ten NIPs remove the eight
rows, the generalization, and the old current formula~0.

Define a parameter reset that preserves the current formula and stack:
\[
 \mathrm{RESET}_a
  =[\mathrm{AX}_0]\;\mathbin{+\!+}\;
   \mathrm{NIP}^{10}\;\mathbin{+\!+}\;
   [\mathrm{DROP},\mathrm{NIP}].
\]
For every $(t,w,d,0)$ it either halts (when the restored cache is~1) or
ends in $(t,w,0,0)$.  After the ten NIPs the current formula is~0 and
the cache is $t$.  DROP restores $t$ while pushing~0, and the final NIP
either recognizes $t=1$ or pops that~0 into the cache.

Use $\mathrm{DISCARD}_a$ in every axiom-leaf block and append
$\mathrm{RESET}_a$ after the ordinary selection of the desired emitted
row or bundle.  The selected formula, underlying stack, zero parameter,
and zero cache then satisfy exactly the old block contract (with early
halting allowed).  The old GEN construction already calls AX with
$a=c$, so its residual is zero.  The successful MP construction also
uses $a=c$ in its following AX reset.  Those two constructions therefore
continue to satisfy their old contracts once their premise blocks have
been replaced.

For completeness, here are all block constructors explicitly.  In the
display below juxtaposition means concatenation, $N=\mathrm{NIP}$,
$D=\mathrm{DROP}$, $A(c,x,y,u)=\mathrm{AX}(c,x,y,u)$, and $L_X$ denotes
the already compiled premise block with canonical output $(X,w,0,0)$.
The displayed row index is its position in the emitted stack:
$j=1$ is $\rho_1$, $j=6$ is $\rho_6$, $j=7$ is $G$, and $j=8$ is $L$.
The eight logical rows, generalization, and old current formula occupy
the first ten stack positions.
\[
\begin{aligned}
L_{\top}
  &=\mathrm{DISCARD}_a,\\
L_{\operatorname{row}_j(a,b,c,u)}
  &=\mathrm{DISCARD}_a\,[A(c,a,b,u)]\,N^j\,[D]\,N^{11-j}
       \,\mathrm{RESET}_a \quad(1\leq j\leq8),\\
L_{\operatorname{gated\ bundle}(a,b,c)}
  &=\mathrm{DISCARD}_a\,[A(c,a,b,0)]\,N^{10}
       \,\mathrm{RESET}_a,\\
L_{Q_i X}
  &=L_X\,[D,N,A(i,i,0,0)]\,N^9\,[D,N,N],\\
L_B
  &=L_A\,[D]\,L_{I(A,B)}\,[N,D,N,\mathrm{MP}(B),A(B,0,0,0)]
       \,N^{10}\,[D,N,N].
\end{aligned}
\]
Each identity is a partial-correctness contract allowing earlier halting.
For the GEN block the emitted generalization is $Q_iX$, its AX has
$a=c=i$, and exactly nine NIPs select that generalization.  For the MP
block the two premise blocks arrange the cached implication $I(A,B)$
and current antecedent $A$; successful MP sets the parameter to $B$.
The following AX has $a=c=B$, and its ten NIPs restore the desired result
without a residual.  If a tested intermediate cache is code~1, the
machine has already halted on the required refutation.

Induction on a native derivation now compiles it into a block with the
old canonical interface.  An initial $\mathrm{AX}_0$, even from an
arbitrary cache and parameter, clears both of them and has zero residual;
the current formula and enlarged stack may be arbitrary.  A compiled
derivation of code~1 works from that interface and the original final
$[\mathrm{DROP},\mathrm{NIP},\mathrm{NIP}]$ test halts.  The clean reader still strictly decreases every
nonempty proof stream, and the next-proof counter advances by one at
each empty stream.  Thus this fixed finite refutation program is
eventually reached.  The repaired source halts if and only if the
unchanged native calculus derives~1.

Both directions are ordinary arithmetizable arguments in ZFC.  Combining
them with the unchanged native-calculus calibration supplied in the
pinned K219 construction gives, for this repaired \emph{source program},
\[
\mathrm{ZFC}\vdash
 \bigl(\text{the repaired source does not halt}\bigr)
 \ \Longleftrightarrow\ \operatorname{Con}(\mathrm{ZFC}).
\]
This conclusion uses the baseline calibration theorem; the changed
source's soundness, fairness, and proof-block completeness are the new
arguments supplied here.

\subsection{Measured source cost}

The all-asymmetric source has 205 DAG nodes.  The new clear-6 loop adds
one previously absent subtree; its attachment adds one subtree; and
the root transfer's attachment adds one subtree.  The transfer itself
was already present in the source.  The resulting four examined
associations all have 208 DAG nodes, 4741 primitive occurrences, maximum
depth~38 or~39, and maximum zero run~6.  The 223-state primitive-interface
implementation supplied later has a complete symbolic runtime proof;
its required blank-tape startup calibration remains open.

The standalone \texttt{counter\_salt\_blocks.py} checks the comparator
and the altered gate with an independent compositional source-tree
interpreter.  Its JSON receipt separately records bounded source cases
and symbolic checks of every block constructor above, with arbitrary
underlying stack tails and arbitrary natural parameters.  The symbolic
checker retains formula constructors; it rejects an unresolved branch
comparison instead of assuming its outcome.  Premise blocks in the
GEN/MP checks are explicit induction hypotheses.  These checks support
the stated unbounded source and compiler arguments; they do not assert
an independently certified binary runtime for the repaired source.

